\documentclass[11pt]{article}

\title{Picard Groups and Invertible Sheaves on $\mathrm{Spec}(R)$}
\author{UlamLib seed note}
\date{}

\newtheorem{definition}{Definition}[section]
\newtheorem{lemma}[definition]{Lemma}
\newtheorem{theorem}[definition]{Theorem}
\newtheorem{corollary}[definition]{Corollary}

\begin{document}
\maketitle

\section{Scope}

This note is the flagship advanced algebraic-geometry target in UlamLib. The
intended Lean destination is
\texttt{ULAM/AlgebraicGeometry/PicardSpec.lean}.

\section{Core definitions}

\begin{definition}
An $R$-module $L$ is invertible if there exists an $R$-module $M$ such that
$L \otimes_R M \cong R$.
\end{definition}

\begin{definition}
An invertible sheaf on $\mathrm{Spec}(R)$ is a locally free sheaf of rank one.
\end{definition}

\begin{definition}
The Picard group $\mathrm{Pic}(R)$ is the group of isomorphism classes of
invertible $R$-modules under tensor product.
\end{definition}

\section{Seed theorem cluster}

\begin{theorem}[Local criterion]
An $R$-module is invertible if and only if it is locally free of rank one at
every prime ideal of $R$.
\end{theorem}

\begin{theorem}
Every invertible $R$-module determines an invertible sheaf on
$\mathrm{Spec}(R)$.
\end{theorem}

\begin{theorem}
The assignment from invertible modules to invertible sheaves induces an
identification of $\mathrm{Pic}(R)$ with the group of isomorphism classes of
invertible sheaves on $\mathrm{Spec}(R)$.
\end{theorem}

\section{Formalization backlog}

\begin{enumerate}
\item Prove the local rank-one criterion for invertible modules.
\item Construct the sheaf attached to an invertible module.
\item Extend the bridge to locally free modules of constant rank.
\item Identify $\mathrm{Pic}(R)$ with invertible sheaves on $\mathrm{Spec}(R)$.
\item Add the cohomological formulation when the library support is ready.
\end{enumerate}

\end{document}
